\documentclass[UTF8,11pt,a4paper,fontset=fandol]{ctexart}

\usepackage[margin=1in]{geometry}
\usepackage{amsmath,amssymb,mathtools}
\usepackage{booktabs}
\usepackage{graphicx}
\usepackage{hyperref}
\usepackage[nameinlink,capitalize]{cleveref}
\usepackage[
  backend=biber,
  style=gb7714-2015,
  gbnamefmt=lowercase
]{biblatex}

\addbibresource{references.bib}

\title{科研论文 LaTeX 模板}
\author{Your Name}
\date{\today}

\begin{document}

\maketitle

\begin{abstract}
这是一个可直接用于科研写作的本地 LaTeX 模板。默认使用 XeLaTeX 和
Biber，适合中文、英文或中英混排论文。
\end{abstract}

\section{Introduction}

LaTeX is still the standard toolchain for papers with equations, references,
figures, and reproducible builds \cite{lamport1986latex}. For machine learning
papers, citation and equation-heavy drafts are easier to maintain in a Git-backed
LaTeX project \cite{vaswani2017attention}.

\section{Method}

Use equations directly:
\begin{equation}
  \operatorname{softmax}(x_i) =
  \frac{\exp(x_i)}{\sum_{j=1}^{n}\exp(x_j)}.
\end{equation}

\section{Results}

\begin{table}[htbp]
  \centering
  \caption{Example result table}
  \label{tab:results}
  \begin{tabular}{lrr}
    \toprule
    Method & Accuracy & Cost \\
    \midrule
    Baseline & 87.5 & 1.00 \\
    Proposed & 88.1 & 0.72 \\
    \bottomrule
  \end{tabular}
\end{table}

See \cref{tab:results} for a compact table example.

\printbibliography

\end{document}
